\documentclass[12pt]{article}
\usepackage[utf8]{inputenc}
\usepackage[T1]{fontenc}
\usepackage[english]{babel}
\usepackage{amsmath, amssymb}
\usepackage{geometry}
\geometry{a4paper, margin=2.5cm}

\title{Formalization and Duality of the Dynamic Ridesharing Problem}
\author{Economic Value Maximization Model}
\date{\today}

\begin{document}

\maketitle

\section{Introduction}
This document presents the mathematical formalization of a dynamic ridesharing problem aimed at maximizing the total economic value generated by the system, along with its dual formulation for shadow price analysis.

\section{Parameters and Variables}

\subsection{Sets and Indices}
\begin{itemize}
    \item $T$: Set of discrete time steps $t$.
    \item $P$: Set of passengers $p$.
    \item $D'$: Extended set of drivers $d$ (including ghost vehicles).
\end{itemize}

\subsection{Parameters}
\begin{itemize}
    \item $\Phi(t, p, d)$: Economic value created by matching passenger $p$ with driver $d$ at time $t$.
    \item $q_d$: Total capacity of vehicle $d$ (including the driver).
    \item $t_p^h$: Departure time (availability) of passenger $p$.
    \item $\mathcal{A}$: Set of admissible triplets $(t, p, d)$ satisfying social compatibility constraints ($penalty = 0$).
\end{itemize}

\subsection{Primal Decision Variable}
\begin{itemize}
    \item $x_{t,p,d} \in \{0, 1\}$: Equals 1 if passenger $p$ is transported by driver $d$ at time $t$, 0 otherwise.
\end{itemize}

\section{Primal Problem (Social Maximization)}

The objective is to maximize the total value over the entire simulation duration under capacity, uniqueness, and temporal availability constraints.

\begin{equation*}
\begin{aligned}
& \underset{x}{\text{Maximize}} && Z_P = \sum_{(t,p,d) \in \mathcal{A}} \Phi(t, p, d) \cdot x_{t,p,d} \\
& \text{subject to:} && \\
& (1) \quad \sum_{p \in P} x_{t,p,d} \leq q_d - 1 && \forall t \in T, \forall d \in D' \quad [\lambda_{t,d}] \\
& (2) \quad \sum_{d \in D'} x_{t,p,d} \leq 1 && \forall t \in T, \forall p \in P \quad [\pi_{t,p}] \\
& (3) \quad x_{t,p,d} = 0 && \forall (t,p,d) \text{ s.t. } t < t_p^h \\
& (4) \quad x_{t,p,d} \geq 0 &&
\end{aligned}
\end{equation*}

\section{Dual Problem (Market Prices)}

The dual problem seeks to minimize the total value of system resources (seats and passenger availability).

\begin{equation*}
\begin{aligned}
& \underset{\lambda, \pi}{\text{Minimize}} && Z_D = \sum_{t \in T} \left( \sum_{d \in D'} (q_d - 1) \lambda_{t,d} + \sum_{p \in P} \pi_{t,p} \right) \\
& \text{subject to:} && \\
& \forall (t, p, d) \in \mathcal{A} \text{ with } t \geq t_p^h: && \\
& \lambda_{t,d} + \pi_{t,p} \geq \Phi(t, p, d) \\
& \lambda_{t,d} \geq 0, \quad \pi_{t,p} \geq 0
\end{aligned}
\end{equation*}

\section{Economic Interpretation}

\begin{itemize}
    \item \textbf{Shadow Price of Seats ($\lambda_{t,d}$)}: Represents the scarcity value of a seat in vehicle $d$ at time $t$. If the vehicle is not full, $\lambda_{t,d} = 0$.
    \item \textbf{Passenger Rent ($\pi_{t,p}$)}: Represents the intrinsic value of the passenger to the system at time $t$, after deducting the "rent" for their seat.
    \item \textbf{Complementarity}: At optimality, a passenger is matched ($x_{t,p,d} = 1$) only if the sum of their resource values ($\lambda + \pi$) is exactly equal to the value $\Phi$ they generate.
\end{itemize}

\end{document}
